% ============================================================================
% COURS 02 — FICHES DÉTAILLÉES DES LIAISONS
% Contenu restauré à partir du Word original.
% ============================================================================

\subsection{Fiches détaillées des dix liaisons usuelles}

Pour identifier une liaison, il faut toujours préciser son \textbf{nom} et ses
\textbf{éléments géométriques caractéristiques} : centre, axe, direction ou
normale. Le symbole normalisé ne décrit pas la solution constructive ; il
représente uniquement les mouvements relatifs autorisés.

\newtcolorbox{TLLiaisonCard}[2][]{
  enhanced,breakable,
  colback=white,colframe=TLCoverBlue,
  colbacktitle=TLCoverBlue,coltitle=white,
  title={#2},fonttitle=\bfseries,
  boxrule=.75pt,arc=2mm,
  left=4mm,right=4mm,top=3mm,bottom=3mm,
  before skip=.9\baselineskip,after skip=.9\baselineskip,
  #1
}

\begin{TLLiaisonCard}{Liaison pivot d'axe \((O,\vec z)\)}
\begin{minipage}[c]{.18\linewidth}\centering
\TLSymPivot
\end{minipage}\hfill
\begin{minipage}[c]{.77\linewidth}
\textbf{Mouvement autorisé :} une rotation autour de l'axe \((O,\vec z)\).

\textbf{Mobilité :} \(R_z\), soit un degré de liberté.

\[
\left\{\mathcal V_{2/1}\right\}_O=
\left\{\begin{array}{cc}
0&0\\0&0\\\omega_z&0
\end{array}\right\}_{\mathcal B}.
\]

\textbf{Exemples technologiques :} arbre guidé par deux roulements, charnière,
axe dans deux paliers. La forme du torseur est valable en tout point de l'axe.
\end{minipage}
\end{TLLiaisonCard}

\begin{TLLiaisonCard}{Liaison glissière de direction \(\vec z\)}
\begin{minipage}[c]{.18\linewidth}\centering
\TLSymGlissiere
\end{minipage}\hfill
\begin{minipage}[c]{.77\linewidth}
\textbf{Mouvement autorisé :} une translation rectiligne suivant \(\vec z\).

\textbf{Mobilité :} \(T_z\), soit un degré de liberté.

\[
\left\{\mathcal V_{2/1}\right\}_P=
\left\{\begin{array}{cc}
0&0\\0&0\\0&v_z
\end{array}\right\}_{\mathcal B}.
\]

\textbf{Paramètre de mouvement :} \(\lambda(t)\), avec
\(\vec{OP}=\lambda(t)\vec z\) et \(v_z=\dot\lambda\).
\end{minipage}
\end{TLLiaisonCard}

\begin{TLLiaisonCard}{Liaison pivot glissant d'axe \((O,\vec z)\)}
\begin{minipage}[c]{.18\linewidth}\centering
\TLSymPivotGlissant
\end{minipage}\hfill
\begin{minipage}[c]{.77\linewidth}
\textbf{Mouvements autorisés :} rotation autour de \(\vec z\) et translation
suivant \(\vec z\), indépendantes l'une de l'autre.

\textbf{Mobilités :} \(R_z\) et \(T_z\), soit deux degrés de liberté.

\[
\left\{\mathcal V_{2/1}\right\}_P=
\left\{\begin{array}{cc}
0&0\\0&0\\\omega_z&v_z
\end{array}\right\}_{\mathcal B}.
\]

\textbf{Exemples :} arbre coulissant dans un coussinet, douille à billes ronde.
\end{minipage}
\end{TLLiaisonCard}

\begin{TLLiaisonCard}{Liaison hélicoïdale d'axe \((O,\vec z)\)}
\begin{minipage}[c]{.18\linewidth}\centering
\TLSymHelicoidale
\end{minipage}\hfill
\begin{minipage}[c]{.77\linewidth}
La rotation et la translation sont \textbf{liées} par le pas de l'hélice.
Cette liaison ne possède donc qu'un seul degré de liberté.

Pour un pas géométrique \(h\) :
\[
\lambda=\frac{h}{2\pi}\theta,
\qquad
v_z=\frac{h}{2\pi}\omega_z.
\]

\[
\left\{\mathcal V_{2/1}\right\}_O=
\left\{\begin{array}{cc}
0&0\\0&0\\\omega_z&\dfrac{h}{2\pi}\omega_z
\end{array}\right\}_{\mathcal B}.
\]

\textbf{Exemples :} système vis-écrou, vis à billes, mécanisme de réglage.
\end{minipage}
\end{TLLiaisonCard}

\begin{TLLiaisonCard}{Liaison sphérique de centre \(O\)}
\begin{minipage}[c]{.18\linewidth}\centering
\TLSymSpherique
\end{minipage}\hfill
\begin{minipage}[c]{.77\linewidth}
\textbf{Mouvements autorisés :} trois rotations indépendantes autour du centre
\(O\). Aucune translation du centre n'est possible.

\textbf{Mobilités :} \(R_x,R_y,R_z\), soit trois degrés de liberté.

\[
\left\{\mathcal V_{2/1}\right\}_O=
\left\{\begin{array}{cc}
\omega_x&0\\\omega_y&0\\\omega_z&0
\end{array}\right\}_{\mathcal B}.
\]

\textbf{Exemples :} rotule lisse, embout à rotule, roulement à rotule.
La forme ci-dessus n'est valable qu'au centre de la liaison.
\end{minipage}
\end{TLLiaisonCard}

\begin{TLLiaisonCard}{Liaison sphérique à doigt de centre \(O\)}
\begin{minipage}[c]{.18\linewidth}\centering
\TLSymSpheriqueDoigt
\end{minipage}\hfill
\begin{minipage}[c]{.77\linewidth}
Une rotation de la liaison sphérique est supprimée par un doigt. Avec la
rotation autour de \(\vec z\) interdite :

\textbf{Mobilités :} \(R_x,R_y\), soit deux degrés de liberté.

\[
\left\{\mathcal V_{2/1}\right\}_O=
\left\{\begin{array}{cc}
\omega_x&0\\\omega_y&0\\0&0
\end{array}\right\}_{\mathcal B}.
\]

Le doigt doit être défini par sa direction ; une autre orientation conduit à
une autre composante de rotation annulée.
\end{minipage}
\end{TLLiaisonCard}

\begin{TLLiaisonCard}{Liaison appui-plan de normale \(\vec z\)}
\begin{minipage}[c]{.18\linewidth}\centering
\TLSymAppuiPlan
\end{minipage}\hfill
\begin{minipage}[c]{.77\linewidth}
Le contact plan sur plan interdit la translation suivant la normale et les
rotations qui modifieraient le parallélisme des plans.

\textbf{Mobilités :} \(T_x,T_y,R_z\), soit trois degrés de liberté.

\[
\left\{\mathcal V_{2/1}\right\}_P=
\left\{\begin{array}{cc}
0&v_x\\0&v_y\\\omega_z&0
\end{array}\right\}_{\mathcal B}.
\]

La forme est valable en tout point du plan de contact.
\end{minipage}
\end{TLLiaisonCard}

\begin{TLLiaisonCard}{Liaison sphère-plan de normale \(\vec z\)}
\begin{minipage}[c]{.18\linewidth}\centering
\TLSymSpherePlan
\end{minipage}\hfill
\begin{minipage}[c]{.77\linewidth}
Le contact ponctuel interdit uniquement la translation suivant la normale au
plan au point de contact \(P\).

\textbf{Mobilités :} \(T_x,T_y,R_x,R_y,R_z\), soit cinq degrés de liberté.

\[
\left\{\mathcal V_{2/1}\right\}_P=
\left\{\begin{array}{cc}
\omega_x&v_x\\\omega_y&v_y\\\omega_z&0
\end{array}\right\}_{\mathcal B}.
\]

\textbf{Autre désignation :} liaison ponctuelle de normale \(\vec z\).
\end{minipage}
\end{TLLiaisonCard}

\begin{TLLiaisonCard}{Liaison cylindre-plan d'axe \(\vec x\), normale \(\vec z\)}
\begin{minipage}[c]{.18\linewidth}\centering
\TLSymCylindrePlan
\end{minipage}\hfill
\begin{minipage}[c]{.77\linewidth}
Le contact linéique rectiligne interdit la translation suivant la normale et
la rotation qui ferait perdre le parallélisme entre la ligne de contact et le
plan.

\textbf{Mobilités :} \(T_x,T_y,R_x,R_z\), soit quatre degrés de liberté.

\[
\left\{\mathcal V_{2/1}\right\}_P=
\left\{\begin{array}{cc}
\omega_x&v_x\\0&v_y\\\omega_z&0
\end{array}\right\}_{\mathcal B}.
\]

\textbf{Autre désignation :} liaison linéaire rectiligne.
\end{minipage}
\end{TLLiaisonCard}

\begin{TLLiaisonCard}{Liaison sphère-cylindre d'axe \((O,\vec z)\)}
\begin{minipage}[c]{.18\linewidth}\centering
\TLSymSphereCylindre
\end{minipage}\hfill
\begin{minipage}[c]{.77\linewidth}
Le contact linéique annulaire conserve les trois rotations et la translation
suivant l'axe du cylindre.

\textbf{Mobilités :} \(R_x,R_y,R_z,T_z\), soit quatre degrés de liberté.

\[
\left\{\mathcal V_{2/1}\right\}_O=
\left\{\begin{array}{cc}
\omega_x&0\\\omega_y&0\\\omega_z&v_z
\end{array}\right\}_{\mathcal B}.
\]

\textbf{Autre désignation :} liaison linéaire annulaire. La forme est valable
au centre du cercle de contact.
\end{minipage}
\end{TLLiaisonCard}

\begin{TLRemark}[Lecture rapide d'un torseur cinématique]
Chaque composante non nulle de la colonne de gauche correspond à une rotation
autorisée ; chaque composante non nulle de la colonne de droite correspond à
une translation autorisée au point d'expression. Le nombre de paramètres
indépendants donne le nombre de degrés de liberté.
\end{TLRemark}

\subsection{Du contact réel à la liaison modélisée}

\begin{center}
\renewcommand{\arraystretch}{1.25}
\begin{tabularx}{\linewidth}{|>{\bfseries}p{.27\linewidth}|X|p{.22\linewidth}|}
\hline
\rowcolor{TLTableHeader}
Géométrie du contact & Liaison usuelle associée & Nombre de ddl\\ \hline
Plan sur plan & Appui-plan & 3\\ \hline
Sphère sur plan & Sphère-plan ou ponctuelle & 5\\ \hline
Cylindre sur plan & Cylindre-plan ou linéaire rectiligne & 4\\ \hline
Sphère dans cylindre & Sphère-cylindre ou linéaire annulaire & 4\\ \hline
Sphère dans sphère & Sphérique & 3\\ \hline
Cylindre dans cylindre coaxial & Pivot glissant & 2\\ \hline
Association de contacts & Pivot, glissière, hélicoïdale ou sphérique à doigt & 1 ou 2\\ \hline
\end{tabularx}
\end{center}

\begin{TLQuoteBox}
Le nom d'une liaison seul est insuffisant. Écrire par exemple
« pivot d'axe \((A,\vec x)\) » ou « appui-plan de normale \(\vec z\) ».
\end{TLQuoteBox}
